\scanpage{126}
% archive sheet number: 125
% source lines 1--9
Voici quelques \unsure{cogitations|cohitations} dans ce sens, lorsque $C$ est muni d'un foncteur additif exact fidèle\deleted{\unsure{F}} $k$-linéaire $F:C\handwritten{\to}\unsure{\mathbb{M}|M}\mathrm{odf}(k)$, et \unsure{correspond} à ce titre à une algèbre profinie (ass un.) $U$, resp. une coalgèbre ass. unit. $A$, sur le corps $k$ (cf nº6). On \deleted{voit que la donnée d} suppose donné une opération $\otimes$, compatible avec $F$, i.e. un\deleted{\unsure{e}} hom. diagonal $U\handwritten{\to}U\overset{\handwritten{\scriptscriptstyle\wedge}}{\otimes}U$ i.e. une composition $A\otimes_kA\handwritten{\to}A$. On se propose de déterminer\inserted{\uline{toutes}} les données de symétrie $\handwritten{\psi}$ pour $\otimes$ (\uline{pas seulement} celles compatibles avec le foncteur $F$ et la donnée de symmétrie ordinaire sur $\mathrm{Modf}(k)$ !) On trouve qu'\unsure{illes|elles} correspondent aux
\margin{\handwritten{[Dans le cas du nº 6, \uline{\unsure{gr. deux|gr. deu.}}, on aura $\sigma=\pi^+\otimes\pi^++\pi^+\otimes\pi^-+\pi^-\otimes\pi^+-\pi^-\otimes\pi^-$]}}
\[
\handwritten{\sigma\in}(U\overset{\handwritten{\scriptscriptstyle\wedge}}{\otimes}U)^{\handwritten{*}}
\tag*{\handwritten{(0)}}
\]
satisfaisant les conditions
\[
\left\{\begin{aligned}
&\handwritten{\sigma'\sigma}=1\quad\text{(i.e. }\handwritten{\sigma}\text{ et }\handwritten{\sigma'}\text{ inverses l'un de l'autre: on aura aussi}\\
&\qquad\handwritten{\sigma\sigma'}=1,\quad\text{puisque }\handwritten{\sigma}\text{ est supposé inversible)}\\
&\boxed{\handwritten{\sigma(\Delta}u)\handwritten{\sigma}^{-1}=\handwritten{(\Delta}u\handwritten{)'}\quad\text{pour tout }u\handwritten{\in}U}
\end{aligned}\right.
\tag*{\handwritten{(1)}}
\]
\margin{\handwritten{condition de symétrie ``tordue'' de l'application diagonale.}}
% source lines 14--18
en désignant par $x\handwritten{\mapsto}x'$ l'automorphisme de symétrie de $U\overset{\handwritten{\scriptscriptstyle\wedge}}{\otimes}U$. NB On ne dit pas que $\handwritten{\Delta}U$ est stable sous $x\handwritten{\leftrightarrow}x'$, mais on considère $\handwritten{\Delta}$ et le symétrique $\handwritten{\Delta'}$ défini par $\handwritten{\Delta'}(u)=\handwritten{\Delta}(u)\handwritten{'}$, qui définit la loi $\otimes$ \uline{opposée} de celle définie par $\handwritten{\Delta}$, et les relations entre $\handwritten{\sigma}$ et $\handwritten{\sigma'}$ sont symétriques et s'expriment\inserted{\handwritten{en disant}} que $\handwritten{\sigma}$ transforme $\handwritten{\Delta}$ en $\handwritten{\Delta'}$, $\handwritten{\sigma'}$ transforme $\handwritten{\Delta'}$ en $\handwritten{\Delta}$.

% source lines 19--22
Lorsqu'il y a une unité bilatère pour $\otimes$ et que $F$ est compatible avec celle-ci,\inserted{i.e. $A$ admet une unité bilatère,} alors la compatibilité de $\handwritten{\psi}$ à l'unité de $C$ s'exprime par les relations
\[
(\mathrm{id}_U\otimes\handwritten{\varepsilon}_U)(\handwritten{\sigma})=(\handwritten{\varepsilon}_U\otimes\mathrm{id}_U)(\handwritten{\sigma})=1,\quad\text{i.e.}\quad\handwritten{\varepsilon\in}1+U^{\handwritten{+}}\overset{\handwritten{\scriptscriptstyle\wedge}}{\otimes}U^{\handwritten{+}}.
\tag*{\handwritten{(2)}}
\]

% source lines 23--26 and bottom handwritten addition
Lorsqu'on a une donnée d'associativité \deleted{sur} $\handwritten{\phi}$ pour $\unsure{\otimes|\varpi}$, compatible avec $F$ et l'associativité ordinaire de $\otimes$ sur $\mathrm{Modf}(k)$, i.e. $A$ est associative, alors la compatibilité entre $\handwritten{\phi}$ et $\handwritten{\psi}$ s'exprime par une relation du style
\[
\handwritten{\deleted{\unsure{\sigma}_1\unsure{\sigma}_2\unsure{\sigma}_3\unsure{\sigma'}_1\unsure{\sigma'}_2\unsure{\sigma'}_3=1\unsure{\sigma}}\quad(\mathrm{id}_U\widehat{\otimes}\Delta)(\sigma)\deleted{\unsure{\sigma^{23}\sigma^{12}|\sigma^{12}\sigma^{13}}}=\deleted{1}\sigma^{12}\sigma^{13}}
\tag*{\handwritten{(3)}}
\]
\handwritten{où les $\sigma^{ij}$\inserted{\unsure{divers}} sont les 3 images de $\sigma$ dans $U\otimes U\otimes U$ par $x\otimes y\mapsto x\otimes y\otimes1$ et les deux autres appl. analogues.}
